\clearpage
\section{总结与展望}

\subsection{项目完成情况}
本项目完成了由 Raft 键值集群、独立 MVCC 事务服务、一致性验证和观测控制台组成的分布式数据库系统。Raft 用任期、日志索引和配置多数组织选举、复制与读写；事务服务以本地 WAL 确定提交和恢复边界，并提供按日志位置读取历史版本的接口。两套服务分别维护各自的状态和提交顺序。

运行记录把客户端结果和内部状态连接起来。节点切换可以从故障信号追踪到新任期和已确认值的读取；事务恢复可以沿 COMMIT 和 CLR 找到版本发布和撤销位置；一致性验证把调用历史和版本依赖转化为合法顺序、依赖环或业务违规位置。控制台以统一入口呈现这些记录，便于分析一次操作从提交到观测的过程。

\subsection{运行约束与演进方向}
固定五节点配置中，写入确认和默认读屏障都依赖至少三个投票成员；成员变更期间旧、新配置分别计算多数。现有批量写用固定窗口兼顾等待时间与本地持久化开销。下一步可分别观测排队、复制和提交耗时，再依据不同负载下的延迟分布调整批量参数。

事务服务启动时读取完整 WAL，再根据检查点和 SQLite 应用游标定位分析与重做。历史查询依赖保留的版本与提交位置。随着持续运行时间增加，可按历史保留策略划定安全回收水位，协调 WAL 归档、版本回收和检查点推进。
